\documentclass[11pt]{article}
\usepackage[margin=2.3cm]{geometry}
\usepackage[T1]{fontenc}
\usepackage{lmodern}
\usepackage{microtype}
\usepackage{booktabs}
\usepackage{graphicx}
\usepackage[numbers,sort&compress]{natbib}
\usepackage[hidelinks]{hyperref}
\usepackage{xcolor}

\newcommand{\nModels}{24}
\newcommand{\nApi}{15}
\newcommand{\nMakers}{11}
\newcommand{\nVision}{8}
\newcommand{\nLocal}{9}
\newcommand{\localA}{d1-3B}
\newcommand{\localAIdx}{78.4}
\newcommand{\localARank}{10}
\newcommand{\localB}{d3 nano}
\newcommand{\localBIdx}{78.0}
\newcommand{\localBRank}{11}
\newcommand{\localCI}{1.4}
\newcommand{\aboveLocal}{Jev 1.13}
\newcommand{\aboveLocalIdx}{79.0}
\newcommand{\aboveLocalGap}{1.0}
\newcommand{\eceSecond}{d1-3B}
\newcommand{\eceSecondVal}{2.4}
\newcommand{\eceSecondRank}{10}
\newcommand{\omniEnGap}{5.8}
\newcommand{\omniArGap}{30.5}
\newcommand{\kevSmallGap}{16.1}
\newcommand{\dThreeEdgeLite}{9.0}
\newcommand{\dThreeLiteNano}{3.1}
\newcommand{\topText}{Decider V1.1 27B}
\newcommand{\topTextIdx}{83.4}
\newcommand{\secondText}{d1}
\newcommand{\secondTextIdx}{82.7}
\newcommand{\topTextGap}{0.7}
\newcommand{\topTextCI}{1.2}
\newcommand{\topFiveSpread}{1.8}
\newcommand{\topVision}{Jev 1.14}
\newcommand{\topVisionIdx}{85.1}
\newcommand{\topTextVisionRank}{3}
\newcommand{\topVisionTextRank}{8}
\newcommand{\eceMin}{2.2}
\newcommand{\eceMax}{20.3}
\newcommand{\enGapWorst}{7.5}
\newcommand{\arGapBest}{12.5}
\newcommand{\arGapWorst}{23.5}
\newcommand{\boxRight}{78.5}
\newcommand{\boxConf}{48.2}
\newcommand{\boxAuroc}{0.80}
\newcommand{\boxIou}{0.70}
\newcommand{\nCases}{3{,}250}
\newcommand{\nDecisions}{4{,}850}


% DRAFT: remove this line and the \draftnote calls before submitting.
\newcommand{\draftnote}[1]{\textcolor{orange}{[#1]}}

\title{System One Bench: An Independent, Frozen-Data Benchmark for Decision Models}
\author{Ali Alhashimi\\ System One Bench\\ \texttt{ali@systemonebench.dev}}
\date{October 2026}

\begin{document}
\maketitle

\begin{abstract}
Decision models read a situation and answer typed questions---yes or no, one of several options, or an ordered
score---with a probability for every answer, in a single fast call. Within weeks of the first release, more than ten
makers shipped one, and each reports its own results on its own data. We present System One Bench, an independent
benchmark that runs every model on the same frozen data with byte-identical requests: \nCases{} text cases
(\nDecisions{} decisions) across 16 labelled datasets in English, French and Arabic, and 1{,}200 image cases across
6 datasets. Every dataset and image is pinned by a SHA-256 hash, and the runner refuses data that does not match, so
a model added later sees exactly the bytes the first one did. We score accuracy, calibration, latency and cost, and
report the \nModels{} models evaluated so far, \nLocal{} of them open-weight models run locally. The top of the text
board is statistically tied: \topText{} leads at \topTextIdx{}, \topTextGap{} points ahead of \secondText{}, within a
95\% interval of $\pm$\topTextCI{}. Small open weights come close: \localA{} and \localB{}, both under 3~GB and run on
a four-core CPU, place \localARank{}th and \localBRank{}th, within \aboveLocalGap{} points of \aboveLocal{}. Text and
image rankings disagree: the text leader is \topTextVisionRank{}rd on images, and the image leader, \topVision{}, is
\topVisionTextRank{}th on text. Calibration error ranges from \eceMin{}\% to \eceMax{}\% and is largely independent
of accuracy. Reading text from a picture costs at most \enGapWorst{} points in English but \arGapBest{} to
\arGapWorst{} points in Arabic on the same items. Code, data, rules and per-model results are public at
\url{https://systemonebench.dev}.
\end{abstract}

\section{Introduction}

Many decisions inside software are small and repeated: route this ticket, flag this message, pick the next action
for an agent. Large language models can make them, but they generate text that has to be parsed, they rarely say how
sure they are, and each call costs time and tokens. \emph{Decision models}---also called System One models, after
the fast, automatic mode of thinking in \citet{kahneman2011thinking}---take a different shape. A request carries a
\texttt{state} (text, JSON or an image) and a set of named questions, each with a type: \texttt{noul} (yes or no),
\texttt{choice} (one of a list of options), \texttt{score} (one of an ordered set of levels) or \texttt{box} (a
bounding box in an image). The model returns, for every question, a probability distribution over the allowed
answers.

Over September and October 2026, TypeSafe, OpenAI, Perplexity, Liquid AI, Microsoft, Cloudflare, Upstage, Inception,
Together AI, Nace AI, vLLM Semantic Router and independent developers released decision models, most of them through
OpenRouter's Decisions API and several as open weights. Each release reports results chosen and run by its maker, often on a leaderboard the maker or its
community controls. Those numbers do not compare: different data, different prompts, different ways of counting
failures, and latency measured from different places.

System One Bench measures every model the same way. Our contributions are:
\begin{itemize}
  \item \textbf{A frozen benchmark.} \nCases{} text cases across 16 labelled public datasets in three languages,
  and 1{,}200 image cases, each pinned to an exact source revision and checked by hash before any run
  (Section~\ref{sec:data}).
  \item \textbf{A strict protocol.} Byte-identical requests for every model, failures counted as wrong, no tuning on
  test data, latency timed from one machine, and no vendor-reported numbers (Section~\ref{sec:protocol}).
  \item \textbf{Results for \nModels{} models} on accuracy, calibration, latency and cost, with error bars,
  including \nLocal{} open-weight models run locally, and a vision track for the \nVision{} that read images (Section~\ref{sec:results}).
  \item \textbf{Findings} that change how these models should be compared: the tied top, small open-weight models close
  behind it, the split between text and vision, calibration as an independent axis, and the cost of reading Arabic from pixels.
\end{itemize}

\section{Data}\label{sec:data}

Table~\ref{tab:data} lists the datasets. The text track samples 200 cases from each dataset's evaluation split (test, validation or
dev; never train), 150 for each XNLI language, with a fixed seed. Most datasets ask one question per case. Typed decisions
\citep{typeddecisions} asks five questions about one agent or product state, so it holds 400 cases and 2{,}000
decisions. In total the text track has \nCases{} cases and \nDecisions{} decisions.

The English datasets cover topic classification (AG News \citep{zhang2015character}), sentiment (SST-2 and SST-5
\citep{socher2013recursive}), emotion \citep{saravia2018carer}, banking intents with 77 options
\citep{casanueva2020efficient}, reading comprehension (BoolQ \citep{clark2019boolq}) and natural language inference
(XNLI \citep{conneau2018xnli}). XNLI also supplies French and Arabic. The other Arabic datasets cover assistant
scenarios and intents (MASSIVE \citep{fitzgerald2023massive}), offensive language (OSACT4 \citep{mubarak2020osact4}),
tweet sentiment \citep{barbieri2022xlmt}, and passage relevance built from MIRACL \citep{zhang2023miracl}, both as
picking the relevant passage out of eight and as a yes-or-no question about one passage.

The vision track has six datasets of 200 cases. Two reuse the exact SST-2 reviews and Arabic tweets of the text
track, rendered as images, so the difference between the two is what reading from pixels costs. The others are
Oxford-IIIT Pet breeds \citep{parkhi2012cats}, yes-or-no questions from VQAv2 \citep{goyal2017making}, counting from
CountBench \citep{paiss2023teaching} as a nine-level score, and RefCOCO referring expressions \citep{yu2016modeling}
as Box questions. Images are resized so every side is between 280 and 1{,}568 pixels, within every model's limits.

\paragraph{Freezing.} Every source is loaded at a fixed commit. A lock file holds a SHA-256 hash of every case---its
state, questions and labels---and of every image it points to. The runner and the scorer refuse data whose hash does
not match. The 400 rendered-text images are committed to the repository, because two versions of the imaging library
draw different pixels from the same text. A scheduled job rebuilds the whole benchmark from its sources every week
and fails if anything no longer matches. Sources do move: during the project, the owner of two of our image datasets
renamed itself, which the pinned revisions and hashes absorb. Changing any
case, prompt or label is a new benchmark version, and every model is rerun on it.

\begin{table}[t]\centering\small
\caption{Datasets. $n$ is the number of cases; typed decisions asks 5 questions per case.}\label{tab:data}
\begin{tabular}{llrp{8cm}}\toprule Dataset & Language & $n$ & Question \\ \midrule
AG News & English & 200 & News topic, 4 options \\
SST-2 & English & 200 & Movie review sentiment, yes/no \\
SST-5 & English & 200 & Sentiment on a 5-level score \\
DAIR Emotion & English & 200 & Emotion, 6 options \\
Banking77 & English & 200 & Banking intent, 77 options \\
BoolQ & English & 200 & Reading comprehension, yes/no \\
XNLI English & English & 150 & Does one sentence follow from another, 3 options \\
Typed decisions & English & 400 & Agent and product states, 5 questions per case \\
XNLI French & French & 150 & Does one sentence follow from another, 3 options \\
XNLI Arabic & Arabic & 150 & Does one sentence follow from another, 3 options \\
MASSIVE scenario & Arabic & 200 & Assistant request scenario, 18 options \\
MASSIVE intent & Arabic & 200 & Assistant request intent, 20 options \\
OSACT4 offensive & Arabic & 200 & Offensive tweet, yes/no \\
Tweet sentiment & Arabic & 200 & Sentiment, 3 options \\
MIRACL rerank & Arabic & 200 & Pick the relevant passage out of 8 \\
MIRACL pairwise & Arabic & 200 & Is this passage relevant, yes/no \\
\midrule
SST-2, as an image & image & 200 & The 200 SST-2 reviews rendered as a picture, yes/no \\
Arabic sentiment, as an image & image & 200 & The 200 Arabic tweets rendered as a picture, 3 options \\
Oxford-IIIT Pet & image & 200 & Cat or dog breed from a photo, 37 options \\
VQAv2 yes/no & image & 200 & A yes/no question about a photo \\
CountBench & image & 200 & How many objects, 2 to 10, as a 9-level score \\
RefCOCO & image & 200 & Box the person or object a phrase describes, 100 people and 100 objects \\
\bottomrule\end{tabular}\end{table}

\begin{table*}[t]\centering\small
\caption{Text track, ranked by the System One Index (mean accuracy over the 16 datasets, equal weight). $\pm$ is a 95\% interval from binomial sampling error. Calibration is top-label ECE. Latency is the median per call through OpenRouter, timed from one machine; Laya ran locally. Failed answers count as wrong. Best value per column in bold.}\label{tab:text}
\resizebox{\textwidth}{!}{\begin{tabular}{rllrrrrrrr}\toprule \# & Model & Maker & Index & English & Arabic & ECE (\%) & Median ms & \$ / 1k calls & Failed \\ \midrule
1 & Decider V1.1 27B & Perplexity & \textbf{83.4} $\pm$1.2 & 81.1 & \textbf{85.0} & 5.7 & 263 & \$0.009 & 0 \\
2 & d1 & Liquid AI & 82.7 $\pm$1.3 & 81.3 & 83.9 & 2.9 & 188 & \$0.015 & 0 \\
3 & Mercury Decide & Inception & 82.4 $\pm$1.3 & \textbf{83.7} & 80.8 & 13.1 & 276 & free & 0 \\
4 & Microsoft-Decision-1 & Microsoft & 82.0 $\pm$1.3 & 81.0 & 83.2 & 2.8 & 240 & \$0.010 & 0 \\
5 & Clef & Cloudflare & 81.6 $\pm$1.3 & 80.6 & 83.0 & \textbf{2.2} & 310 & \$0.116 & 0 \\
6 & Clef Flash & Cloudflare & 80.3 $\pm$1.3 & 80.2 & 80.5 & 3.5 & 267 & \$0.044 & 0 \\
7 & Clef Omni & Cloudflare & 79.7 $\pm$1.3 & 78.6 & 80.5 & 3.4 & 208 & \$0.074 & 0 \\
8 & Jev 1.14 & TypeSafe & 79.3 $\pm$1.4 & 78.8 & 80.3 & 6.0 & \textbf{136} & \$0.028 & 0 \\
9 & Jev 1.13 & TypeSafe & 79.0 $\pm$1.4 & 78.5 & 79.9 & 6.1 & 150 & \$0.028 & 0 \\
10 & Drex v1.5 & Nace AI & 77.6 $\pm$1.4 & 77.6 & 77.0 & 10.8 & 198 & \$0.010 & 0 \\
11 & Kev 4B & Jared Palmer & 76.8 $\pm$1.4 & 78.3 & 74.8 & 5.1 & 552 & \$0.010 & 0 \\
12 & GPT-6 Luna Decisions & OpenAI & 76.0 $\pm$1.4 & 75.2 & 78.8 & 9.8 & 215 & \$0.040 & 2 \\
13 & Tev1 4B Experimental & Together AI & 71.5 $\pm$1.4 & 66.4 & 75.8 & 5.3 & 232 & \$0.019 & 200 \\
14 & Solar Decide & Upstage & 68.8 $\pm$1.5 & 64.7 & 73.7 & 20.3 & 715 & \$0.042 & 200 \\
15 & Solar Decide Flash & Upstage & 63.6 $\pm$1.5 & 61.3 & 66.5 & 14.4 & 586 & \$0.042 & 200 \\
16 & Laya Router & Convai Innovations & 60.4 $\pm$1.5 & 64.7 & 53.2 & 20.0 & 75$^\dagger$ & local & 0 \\
\bottomrule\end{tabular}}
\\[2pt]\footnotesize $^\dagger$Local, one call at a time on an Apple M5 GPU; not comparable with API latency.
\end{table*}

\begin{table*}[t]\centering\small
\caption{Vision track, ranked by the Vision Index (mean accuracy over the 5 image datasets). Text rank is the model's place on the System One Index in Table~\ref{tab:text}.}\label{tab:vision}
\resizebox{\textwidth}{!}{\begin{tabular}{rlrrrrrrrr}\toprule \# & Model & Vision Index & SST-2 & Arabic sentiment & Oxford-IIIT Pet & VQAv2 yes/no & CountBench & ECE (\%) & Text rank \\ \midrule
1 & Jev 1.14 & 85.1 & 93.5 & 63.5 & 88.5 & 86.5 & 93.5 & 2.9 & 8 \\
2 & Jev 1.13 & 83.9 & 94.5 & 59.0 & 83.5 & 88.0 & 94.5 & 2.8 & 9 \\
3 & Decider V1.1 27B & 83.0 & 96.0 & 51.0 & 84.0 & 89.5 & 94.5 & 5.6 & 1 \\
4 & GPT-6 Luna Decisions & 81.6 & 88.5 & 67.0 & 87.0 & 80.5 & 85.0 & 5.3 & 12 \\
5 & Mercury Decide & 79.4 & 93.0 & 39.0 & 84.5 & 87.0 & 93.5 & 9.9 & 3 \\
6 & Clef & 78.1 & 95.5 & 35.0 & 82.5 & 87.5 & 90.0 & 3.6 & 5 \\
7 & Clef Flash & 78.0 & 95.5 & 33.5 & 85.0 & 88.0 & 88.0 & 9.4 & 6 \\
8 & Clef Omni & 72.2 & 89.5 & 28.5 & 88.0 & 87.0 & 68.0 & 20.6 & 7 \\
\bottomrule\end{tabular}}\end{table*}



\section{Protocol and metrics}\label{sec:protocol}

\paragraph{Requests.} Every model receives byte-identical \texttt{state} and \texttt{questions} bodies through the
same request shape: OpenRouter's Decisions API, the maker's own endpoint, or, for open weights, a local server that
answers the same \texttt{/v1/systemone} request (llama.cpp \citep{gerganov2023llamacpp} for GGUF builds, or the
maker's own server). There are no
per-model prompts, option orders or instructions, and nothing is tuned on the test data. Only transient errors (HTTP
408, 429 and 5xx) are retried. A call that fails, or an answer that cannot be read, counts as wrong.

\paragraph{Accuracy and indexes.} A decision is right when the most probable answer equals the label. The
\emph{System One Index} is the mean accuracy over the 16 text datasets, each weighted equally, on a 0--100 scale; the
English and Arabic sub-indexes average the 8 English and 7 Arabic datasets. The \emph{Vision Index} is the same over
the five image datasets with a single correct answer. Box questions are scored separately: a box is right when its
intersection over union with the labelled box is at least 0.5.

\paragraph{Error bars.} We report a 95\% interval for each index from binomial sampling error, treating datasets as
independent and counting typed-decisions cases rather than decisions, since its five questions share a state. These
intervals are for one model; differences between two models on the same items have smaller sampling error, and a
paired test needs per-case predictions, which we plan to publish with each run\draftnote{Publish per-case predictions (plan: Hugging Face dataset), then report paired tests here.}. As a second check, rerunning one model a week later on
rebuilt data moved individual datasets by up to 1.5 points while leaving its index unchanged at one decimal place.

\paragraph{Calibration.} We report top-label expected calibration error \citep{naeini2015obtaining,guo2017calibration}
over ten equal-width bins, pooled over every decision: the average gap between a model's stated confidence and how
often it is right.

\paragraph{Latency and cost.} Latency is the full round trip from the caller, so it depends on the network. Every API
model was therefore timed again from one cloud server on the same day, with eight requests in flight, over all
\nCases{} cases; from the machine behind our earlier runs, the same calls took 1.5 to 3.5 times as long, and the
ranking changed. Open-weight models run locally, one call at a time and with no network: Laya on an Apple M5 GPU, and
the others on a four-core (4-vCPU) cloud server with no GPU, at 8-bit quantisation (Q8\_0). Their latency is reported
but not compared with the API models'. Cost is what the API billed per call; local runs have none.

\section{Results}\label{sec:results}

Table~\ref{tab:text} gives the text track and Table~\ref{tab:vision} the vision track.

\paragraph{The top of the text board is tied.} \topText{} leads the System One Index at \topTextIdx{}, and
\secondText{} follows at \secondTextIdx{}. The gap, \topTextGap{} points, is well within the $\pm$\topTextCI{}
interval of either model, and the top five models lie within \topFiveSpread{} points. A ranking among them should not
be read as a finding; the differences that matter are in which datasets each model is best at, and in calibration,
latency and cost.

\paragraph{Small open-weight models come close.} \localA{} and \localB{}, both under 3~GB and run on a four-core
CPU, place \localARank{}th and \localBRank{}th at \localAIdx{} and \localBIdx{}, within the $\pm$\localCI{} intervals
of \aboveLocal{} (\aboveLocalIdx{}) and above several API models. \localA{} is also the second best calibrated model on
the board (\eceSecondVal{}\%). The smallest models keep much more of their accuracy in English than in Arabic:
d1-omni-600M is \omniEnGap{} points below the full d1 in English and \omniArGap{} points below it in Arabic, and Kev
0.8B, trained for English, loses \kevSmallGap{} points between its English and Arabic sub-indexes. Within one family,
size buys accuracy: d3 gains \dThreeEdgeLite{} points from edge to lite and \dThreeLiteNano{} from lite to nano.

\paragraph{Text and vision rank models differently.} \topVision{} leads the Vision Index at \topVisionIdx{} but is
\topVisionTextRank{}th on text, while \topText{}, the text leader, is \topTextVisionRank{}rd on images. A single
leaderboard that mixes the two would hide this, so we keep two indexes and never combine them.

\paragraph{Calibration is a separate axis.} Calibration error ranges from \eceMin{}\% to \eceMax{}\% across models,
and is only loosely related to accuracy: the model with the best English sub-index has one of the highest
calibration errors, while the best-calibrated model is ranked fifth and the second best, \eceSecond{}, is an open-weight model
ranked \eceSecondRank{}th. For a system that acts on a model's
confidence---escalating uncertain cases to a person, for example---this matters as much as accuracy.

\paragraph{Reading Arabic from pixels is hard.} On the same 200 SST-2 reviews, reading the text from an image costs
every model at most \enGapWorst{} points, and some lose nothing. On the same 200 Arabic tweets, it costs \arGapBest{}
to \arGapWorst{} points. The gap suggests that these models' image encoders recognise Arabic script far less
reliably than Latin script, which a text-only evaluation would not show.

\paragraph{Interface limits become failures.} Two model families accept at most 20 or 26 options per question, so
every Banking77 case (77 options) fails and counts as wrong. We count these as errors rather than skipping the
dataset, because a user who needs 77 options would hit the same limit.

\paragraph{Box questions.} The one model that answers Box questions gets \boxRight{}\% of boxes right at an
intersection over union of 0.5 or more (mean \boxIou{}), while stating an average confidence of \boxConf{}\%. It is
underconfident, but its confidence still ranks answers well (AUROC \boxAuroc{}), so it works as a threshold if not as a
probability.

\section{Related work}

Benchmarks for general language understanding such as GLUE \citep{wang2018glue} established fixed test sets and a
single aggregate score; we follow them in averaging over datasets but add calibration, latency and cost, which matter
for decisions made inside software. Calibration of neural classifiers has been studied since \citet{guo2017calibration}
showed that modern networks are often overconfident. Several community leaderboards track decision models, including
the Jev Decision Index on Hugging Face and JevBench. They differ from System One Bench in data (some partly private),
in how they score abstentions and failures, and in whether they report latency, cost and calibration.
\draftnote{Add citations for these leaderboards once stable URLs are confirmed.}

\section{Limitations}

Our datasets are public, and some may overlap the training data of the models we test; some models disclose
training on the train splits of datasets whose evaluation splits we use (the Laya Arabic fine-tunes on XNLI and
MASSIVE, and the Naluri encoders on XNLI, MASSIVE and DAIR Emotion), and we note this where it is known. Each dataset has 150 to
400 cases, so per-dataset accuracies carry intervals of several points. Latency through an API depends on the route
between caller and provider, so our numbers compare models with each other, not with a user's own measurements; local runs depend on the
hardware and are not compared at all. Open-weight models ran as 8-bit quantised builds, which may score slightly
differently from the original weights. API
models can change behind a stable name, and one model was retired during the project; we record the served model
version where the API reports it and keep retired runs for reference. Finally, the text track
is weighted towards English (8 of 16 datasets) and covers only three languages.

\section{Conclusion}

System One Bench compares decision models under one protocol on frozen data, and reports what makers' own
leaderboards tend to leave out: error bars, calibration, latency from one place, cost, and failures. It is a living
benchmark: new models are added within days of release, open weights run on the same data as the APIs, every
round is a tagged release, and anyone can request a
model or rerun the benchmark with the published code.

\bibliographystyle{plainnat}
\bibliography{refs}

\end{document}
